\documentclass[a4paper, serif, 10pt]{moderncv}

%% moderncv
% color and style
\moderncvcolor{black}
\moderncvstyle{banking}

% renew moderncv command to adapt for CJK colon
\renewcommand*{\cvitem}[3][.25em]{%
  \ifstrempty{#2}{}{\hintstyle{#2}：}{#3}%
  \par\addvspace{#1}}

\renewcommand*{\cvitemwithcomment}[4][.25em]{%
  \savebox{\cvitemwithcommentmainbox}{\ifstrempty{#2}{}{\hintstyle{#2}：}#3}%
  \setlength{\cvitemwithcommentmainlength}{\widthof{\usebox{\cvitemwithcommentmainbox}}}%
  \setlength{\cvitemwithcommentcommentlength}{\maincolumnwidth-\separatorcolumnwidth-\cvitemwithcommentmainlength}%
  \begin{minipage}[t]{\cvitemwithcommentmainlength}\usebox{\cvitemwithcommentmainbox}\end{minipage}%
  \hfill% fill of \separatorcolumnwidth
  \begin{minipage}[t]{\cvitemwithcommentcommentlength}\raggedleft\small\itshape#4\end{minipage}%
  \par\addvspace{#1}}

%% page layout/margins
\usepackage[top=2.5cm, bottom=2.5cm, left=1.5cm, right=1.5cm]{geometry}

\nopagenumbers{}

%% Babel config for English language
\usepackage[english]{babel}

%% fontspec
\usepackage{fontspec}

\IfFontExistsTF{Linux Libertine}{
  \setmainfont[Ligatures={TeX, Common}, Numbers=Lining, ItalicFont=Linux Libertine]{Linux Libertine}
}{}
\IfFontExistsTF{Linux Libertine O}{
  \setmainfont[Ligatures={TeX, Common}, Numbers=Lining, ItalicFont=Linux Libertine O]{Linux Libertine O}
}{}

%% CTeX
% CJK support, used to show CJK characters in the resume
%
% - heading=false: disable ctex heading
% - punct=kaiming: use kaiming punctuations styles for CJK
% - scheme=plain: use plain scheme, do not override \normalsize font size
% - space=auto: space settings for CJK characters
%
% CJK fonts are resolved per locale language, e.g. Noto Serif/Sans CJK JP
% for Japanese resumes, with fallback to the next region if not installed
%
% ref:
% - http://ctan.mirrorcatalogs.com/language/chinese/ctex/ctex.pdf
\usepackage[UTF8, heading=false, punct=kaiming, scheme=plain, space=auto]{ctex}

\IfFontExistsTF{Noto Serif CJK SC}{
  \setCJKmainfont{Noto Serif CJK SC}
}{}
\IfFontExistsTF{Noto Sans CJK SC}{
  \setCJKsansfont{Noto Sans CJK SC}
}{}

%% line spacing
\usepackage{setspace}
\setstretch{1.125}

%% URL styling - use normal text instead of monospace
\urlstyle{same}

%% Patch \httplink and \httpslink to handle full URLs
% The original moderncv macros always prepend http:// or https:// to URLs.
% This causes issues when the URL already has a protocol (e.g., https://example.com)
% resulting in malformed URLs like https://https://example.com.
% This patch checks if the URL already contains "://" and uses it directly if so.
% Note: We use \str_set:Nx (x-expansion) to expand macros like \@homepage first.
\ExplSyntaxOn
\str_new:N \g_yamlresume_url_str

\RenewDocumentCommand{\httpslink}{O{}m}{
  \str_set:Nx \g_yamlresume_url_str {#2}
  \str_if_in:NnTF \g_yamlresume_url_str {://}
    { \url{#2} }
    { \url{https://#2} }
}

\RenewDocumentCommand{\httplink}{O{}m}{
  \str_set:Nx \g_yamlresume_url_str {#2}
  \str_if_in:NnTF \g_yamlresume_url_str {://}
    { \url{#2} }
    { \url{http://#2} }
}
\ExplSyntaxOff

\name{张伟}{}
\title{高级机械工程师 | 汽车行业 | 10年经验}
\phone[mobile]{+86 138 1234 5678}
\email{zhangwei.mech@qq.com}
\homepage{https://www.linkedin.com/in/zhangwei-mech}

\address{中国，上海市，上海，上海市浦东新区世纪大道100号，200120}{}{}

\extrainfo{{\small \faLinkedin}\ \href{https://www.linkedin.com/in/zhangwei-mech}{@zhangwei-mech} {} {} {} • {} {} {} 
{\small \faGithub}\ \href{https://github.com/zhangwei-mech}{@zhangwei-mech}}

\begin{document}

\maketitle

\section{简介}

\cvline{}{\begin{itemize}
\item 机械工程硕士，10年汽车行业机械设计与开发经验，熟悉从概念设计到量产的全流程。
\item 精通SolidWorks、CATIA、AutoCAD等CAD软件，熟练使用ANSYS、Abaqus进行有限元分析。
\item 擅长汽车零部件结构设计、轻量化优化、公差分析及GD\&T，具备丰富的项目管理经验。
\item 具备良好的团队协作和沟通能力，曾领导5人团队完成多个重点项目，获得上海市科技进步二等奖。
\item 持续学习新技术，关注新能源汽车和智能制造领域，致力于推动机械设计的创新与优化。
\end{itemize}}
\section{教育背景}

\cventry{2010年9月–2013年6月}
        {硕士，机械工程，成绩：3.8/4.0}
        {上海交通大学}
        {\url{https://www.sjtu.edu.cn}}
        {}
        {\begin{itemize}
\item 研究方向为汽车轻量化设计与结构优化，参与国家自然科学基金项目。
\item 发表2篇EI论文，获得校级优秀硕士论文。
\item 担任助教，辅导本科生机械设计课程。
\end{itemize}
\textbf{课程}：高等机械设计、有限元分析、热力学与流体力学、机器人学、控制工程、材料科学与工程、先进制造技术、汽车系统动力学}

\cventry{2006年9月–2010年6月}
        {学士，机械设计制造及其自动化，成绩：3.6/4.0}
        {华中科技大学}
        {\url{https://www.hust.edu.cn}}
        {}
        {\begin{itemize}
\item 连续三年获得校级奖学金，担任机械创新协会会长。
\item 参加全国大学生机械创新设计大赛，获湖北省一等奖。
\item 毕业设计为“小型机器人机械臂设计”，获优秀毕业设计。
\end{itemize}
\textbf{课程}：机械原理、机械设计、工程力学、电工电子技术、微机原理与接口技术、工程材料、互换性与技术测量、机械制造技术基础}
\section{工作经历}

\cventry{2018年7月至今}
        {高级机械工程师}
        {上海汽车集团股份有限公司}
        {\url{https://www.saicmotor.com}}
        {}
        {\begin{itemize}
\item 负责新能源汽车电池包结构设计与优化，主导完成3款车型的电池包开发，实现减重15\%并满足碰撞安全要求。
\item 使用CATIA进行3D建模和2D图纸绘制，运用ANSYS进行静力学、模态和疲劳分析，优化结构设计。
\item 制定GD\&T方案，与供应商进行技术沟通，确保零部件制造质量，降低生产成本10\%。
\item 领导5人跨部门团队，协调设计、仿真、工艺和采购，确保项目按时交付。
\item 编写技术文档和设计规范，培训新员工，提升团队整体设计能力。
\end{itemize}
\textbf{关键字}：电池包设计、轻量化、CATIA、ANSYS、GD\&T、项目管理}

\cventry{2013年7月–2018年6月}
        {机械工程师}
        {博世汽车部件（苏州）有限公司}
        {\url{https://www.bosch.com.cn}}
        {}
        {\begin{itemize}
\item 负责动力总成零部件的机械设计，包括变速箱壳体、支架和连接件，年产量超50万件。
\item 使用SolidWorks进行详细设计，完成公差分析和尺寸链计算，确保装配精度。
\item 参与DFMEA和PFMEA分析，识别设计风险并制定改进措施，产品不良率降低20\%。
\item 与德国总部团队协作，引入新设计标准，优化产品结构，缩短开发周期15\%。
\item 获得博世全球创新奖提名，连续两年被评为优秀员工。
\end{itemize}
\textbf{关键字}：动力总成、SolidWorks、公差分析、DFMEA、制造工艺}
\section{语言}

\cvline{中文}{母语或双语水平}
\cvline{英语}{完全专业水平 \hfill \textbf{关键字}：CET-6、TOEFL 105、技术文档写作}
\section{专业技能}

\cvline{CAD软件}{专家 \hfill \textbf{关键字}：SolidWorks、CATIA、AutoCAD、Creo、NX}
\cvline{仿真分析}{高级 \hfill \textbf{关键字}：ANSYS、Abaqus、HyperMesh、MATLAB、Simulink}
\cvline{机械设计}{专家 \hfill \textbf{关键字}：结构设计、轻量化、公差分析、GD\&T、材料选型}
\cvline{项目管理}{中级 \hfill \textbf{关键字}：PMP、Agile、跨部门协作、风险管理}
\cvline{编程与工具}{中级 \hfill \textbf{关键字}：Python、C++、MATLAB、Git}
\section{荣誉}

\cventry{2020年12月}
        {上海市人民政府}
        {上海市科技进步二等奖}
        {}
        {}
        {因在新能源汽车电池包轻量化设计方面的突出贡献，作为主要完成人获奖。}

\cventry{2019年12月}
        {上海汽车集团股份有限公司}
        {上汽集团优秀员工}
        {}
        {}
        {年度绩效排名前5\%，表彰在电池包开发项目中的卓越领导和团队协作。}
\section{证书}

\cventry{2021年6月}
        {上海市人力资源和社会保障局}
        {高级机械工程师职称}
        {}
        {}
        {}

\cventry{2018年9月}
        {项目管理协会（PMI）}
        {PMP认证}
        {}
        {}
        {}

\cventry{2015年3月}
        {Dassault Systèmes}
        {SolidWorks认证专家}
        {}
        {}
        {}
\section{出版刊物}

\cventry{2019年5月}
        {基于有限元分析的汽车悬架结构优化设计}
        {机械工程学报}
        {\url{https://doi.org/10.3901/JME.2019.05.001}}
        {}
        {\begin{itemize}
\item 提出一种基于响应面法的悬架结构优化方法，实现减重12\%并提高疲劳寿命。
\item 通过ANSYS仿真和台架试验验证，结果误差小于5\%。
\item 被引次数超过30次，获期刊年度优秀论文。
\end{itemize}}

\cventry{2017年11月}
        {轻量化材料在汽车车身中的应用研究}
        {汽车工程}
        {\url{https://doi.org/10.19562/j.chinasae.qcgc.2017.11.001}}
        {}
        {\begin{itemize}
\item 对比分析高强钢、铝合金和碳纤维复合材料在车身中的应用效果。
\item 提出多材料混合车身设计方案，实现减重20\%的同时满足碰撞安全法规。
\item 为后续车型开发提供材料选型依据。
\end{itemize}}
\section{项目}

\cventry{2020年1月–2021年12月}
        {针对纯电动车型的电池包进行结构设计、轻量化和安全优化。}
        {新能源汽车电池包结构设计与优化}
        {\url{https://www.saicmotor.com/tech/battery}}
        {}
        {\begin{itemize}
\item 设计电池包上下壳体、模组支架和冷却板，满足IP67防护等级和振动疲劳要求。
\item 采用拓扑优化和尺寸优化方法，实现电池包减重15\%，能量密度提升8\%。
\item 通过CAE仿真和物理试验验证，顺利通过国标GB 38031碰撞安全测试。
\item 项目投产后，年产量达10万台，新增产值超5亿元。
\end{itemize}
\textbf{关键字}：电池包、轻量化、拓扑优化、碰撞安全}

\cventry{2018年3月–2019年12月}
        {对某轿车底盘进行轻量化设计，采用铝合金和液压成型技术。}
        {某车型底盘轻量化项目}
        {}
        {}
        {\begin{itemize}
\item 重新设计前副车架和后扭力梁，采用铝合金材料，减重25\%。
\item 使用HyperMesh进行网格划分和优化，通过Abaqus进行强度刚度分析。
\item 协调供应商完成样件制造和台架试验，满足耐久性要求。
\item 项目成果应用于两款车型，年节约成本约2000万元。
\end{itemize}
\textbf{关键字}：底盘、铝合金、液压成型、轻量化}
\section{兴趣爱好}

\cvline{技术}{3D打印、机器人技术、智能制造}
\cvline{户外}{徒步、摄影、骑行}
\section{志愿服务}

\cventry{2016年1月–2020年12月}
        {技术讲师}
        {中国机械工程学会}
        {\url{https://www.cmes.org}}
        {}
        {\begin{itemize}
\item 作为志愿者讲师，为高校学生和初入职场工程师讲授机械设计实践课程。
\item 组织“机械创新设计”工作坊，累计培训学员超过500人。
\item 参与学会年度学术会议的组织工作，负责分会场协调。
\end{itemize}}
    
\end{document}